%=======================================================================
%  CSE558 Data Science - Midterm Project Report Template
%  Monsoon 2026, IIIT-Delhi
%
%  How to use this template
%   1. Set your project title in the EDIT THIS block below.
%   2. Replace everything in LIGHT BLUE with your own text and numbers.
%      The numbers in the example are made up for illustration.
%   3. Grey italic lines are instructions. Delete them as you write, or
%      hide them all at once with \showhintsfalse.
%   4. Upload your plots to the project and put their file names in
%      \figfile{...}{...}. Until a file exists, a labelled box is shown.
%
%  Sections 1-5 follow items 1-5 of the report guidelines; Section 6 is
%  the limitations of your analysis. The link to your code (item 6) goes
%  at the top. At most 8 pages.
%  Compile with pdfLaTeX (Overleaf's default setting works), twice.
%=======================================================================
\documentclass[10pt,a4paper,twocolumn]{article}

%========================= EDIT THIS =========================
% Your project title: it sets the heading AND the PDF's title.
\newcommand{\projecttitle}{Data Science Midterm Project Report\\[2pt]Team: Graph Out Loud}
\newcommand{\projectsubtitle}{Delayed in New York: does weather, airline or season decide it,
  and can a late flight catch up in the air?}
\newif\ifshowhints
\showhintsfalse     % \showhintsfalse hides every grey instruction line
%=============================================================

%-------------------------------------------- fonts (no need to edit)
\usepackage{iftex}
\ifPDFTeX
  \usepackage[T1]{fontenc}
  \usepackage[utf8]{inputenc}
  \IfFileExists{carlito.sty}%
    {\usepackage[sfdefault]{carlito}}%              Calibri look
    {\usepackage[scaled=0.92]{helvet}\renewcommand{\familydefault}{\sfdefault}}%
  \usepackage{courier}
\else
  \usepackage[no-math]{fontspec}
  \setmainfont{Carlito}[Extension=.ttf, UprightFont=*-Regular, BoldFont=*-Bold,
    ItalicFont=*-Italic, BoldItalicFont=*-BoldItalic]
\fi

%----------------------------------------------------------- packages
\usepackage[margin=18mm,bottom=20mm,columnsep=7mm]{geometry}
\usepackage{amsmath}
\usepackage{xcolor}
\usepackage{graphicx}
\usepackage{tabularx}
\usepackage{array}
\usepackage{booktabs}
\usepackage{enumitem}
\usepackage{titlesec}
\usepackage{needspace}
\usepackage[font=small,labelfont=bf]{caption}
\usepackage[hidelinks]{hyperref}
\hypersetup{pdftitle={Data Science Midterm Project Report - Team Graph Out Loud},
            pdfsubject={CSE558 Data Science, Monsoon 2026: Midterm Project Report}}

%------------------------------------------------------------- layout
\definecolor{eg}{HTML}{000000}         % our own text (was light blue in the template)
\definecolor{hintgrey}{gray}{0.58}     % instructions (light grey)

\titleformat{\section}{\normalfont\large\bfseries}{\thesection.}{0.5em}{}
\titlespacing*{\section}{0pt}{12pt}{3pt}
\titleformat{\subsection}{\normalfont\normalsize\bfseries}{\thesubsection}{0.5em}{}
\titlespacing*{\subsection}{0pt}{10pt}{2pt}
\newcommand{\sectionbreak}{\needspace{10\baselineskip}}
\newcommand{\subsectionbreak}{\needspace{10\baselineskip}}

\setlength{\parindent}{0pt}
\setlength{\parskip}{4pt}
\setlist{leftmargin=1.2em, itemsep=2pt, topsep=2pt, parsep=0pt}
\renewcommand{\arraystretch}{1.25}
\setlength{\tabcolsep}{4pt}
\newcolumntype{L}[1]{>{\raggedright\arraybackslash}p{#1}}
\newcolumntype{Y}{>{\raggedright\arraybackslash}X}
\newcolumntype{K}[1]{>{\raggedright\arraybackslash\bfseries\boldmath}p{#1}}  % bold labels

%--------------------------------------------------------- shorthands
% Instruction line. Delete these as you fill the template in.
\newcommand{\hint}[1]{\ifshowhints{\footnotesize\itshape\color{hintgrey}#1\par}\fi}
% Worked example: inline or in a table cell. REPLACE the contents.
\newcommand{\eg}[1]{\textcolor{eg}{#1}}
% Worked example running over a whole paragraph.
\newcommand{\egblock}[1]{{\color{eg}#1\par}}
% Bold label followed by an example.
\newcommand{\field}[2]{\textbf{#1}~\eg{#2}\par}
% Part a)-c) of a test, and a small bold heading.
\newcommand{\tpart}[2]{\par\needspace{12\baselineskip}\smallskip\textbf{#1)~#2}\par}
\newcommand{\minihead}[1]{\par\needspace{8\baselineskip}\smallskip\textbf{#1}\par}
% Column and file names.
\newcommand{\code}[1]{\texttt{#1}}
% \figfile{file}{height}: your plot, or a labelled box until the file exists.
\newcommand{\figfile}[2]{%
  \IfFileExists{#1}%
    {\includegraphics[width=\linewidth,height=#2,keepaspectratio]{#1}}%
    {\fbox{\parbox[c][#2][c]{\dimexpr\linewidth-2\fboxsep-2\fboxrule\relax}%
      {\centering\small\itshape\color{hintgrey}Your plot: \detokenize{#1}}}}}

%=======================================================================
\begin{document}

\twocolumn[%
{\small\bfseries CSE558 Data Science\ \ \textbullet\ \ Midterm Project Report\ \ \textbullet\ \ Monsoon 2026\par}
\vspace{4pt}
{\LARGE\bfseries\raggedright\projecttitle\par}
\vspace{4pt}
{\large\raggedright\projectsubtitle\par}
\vspace{6pt}
\rule{\linewidth}{1pt}

{\small
\begin{tabularx}{\linewidth}{@{}K{2.4cm} Y@{}}
\toprule
Group name & Graph Out Loud \\
Dataset & BTS On-Time Performance, all 2013 departures from JFK, LGA and EWR (transtats.bts.gov), with NOAA hourly airport weather; packaged as github.com/hadley/nycflights13 and pypi.org/project/nycflights13 \\
Code & \href{https://github.com/Nikkumar-12345/cse558-midterm-graph-out-loud}{github.com/Nikkumar-12345/cse558-midterm-graph-out-loud} (\code{analysis.py} and README) \\
\bottomrule
\end{tabularx}}

\vspace{2pt}
{\small
\begin{tabularx}{\linewidth}{@{}L{0.5cm} L{3.0cm} L{1.5cm} L{4.2cm} Y@{}}
\toprule
\textbf{\#} & \textbf{Name} & \textbf{Roll no.} & \textbf{Institute email} & \textbf{Contribution} \\
\midrule
1 & Vatsal Garg & 2024611 & vatsal24611@iiitd.ac.in & Dataset and preprocessing (1, 2, 3) \\
2 & Nikhil Kumar & 2024382 & nikhil24382@iiitd.ac.in & Exploratory data analysis and problem statements (4) \\
3 & Shakti & 2024519 & shakti24519@iiitd.ac.in & In-flight delay recovery and weather impact analysis (5: P1, P2) \\
4 & Navanshu Nagar & 2024374 & navanshu24374@iiitd.ac.in & Carrier performance and seasonal cancellation analysis (5: P3, P4) \\
5 & M.\ Rayhankhan & 2022269 & rayhan22269@iiitd.ac.in & Weekly delay patterns, synthesis, limitations, code and README (5, 6) \\
\bottomrule
\end{tabularx}}
\vspace{12pt}
]

%-----------------------------------------------------------------------
\section{Dataset}

\field{Source:}{U.S. Bureau of Transportation Statistics (BTS) On-Time Performance records for
every domestic flight that left New York City's three airports (JFK, LaGuardia, Newark) in 2013,
as released in the \code{nycflights13} package. One row is one scheduled flight. We join the
package's \code{weather} table (NOAA hourly observations at the same airports, 26{,}115 rows
$\times$ 15 columns) on airport and hour.}

\field{What is raw about it:}{clock times are integers (\code{517} is 05:17, and \code{2400}
appears 179 times); a cancelled flight is not labelled and shows up only as a missing
\code{dep\_time}; a diverted flight shows up as a missing \code{arr\_delay}; the weather table has
a wind speed of 1{,}048 mph, 165 missing hours and a gust column that is 79.6\% empty; delays run
from $-86$ to $+1{,}301$ minutes.}

{\small
\begin{tabularx}{\linewidth}{@{}L{3.0cm} Y Y@{}}
\toprule
 & \textbf{Raw file} & \textbf{After encoding} \\
\midrule
Instances (rows) & 336{,}776 & 336{,}776 \\
Attributes (columns) & 19 & 167 \\
Instances $\times$ attributes & 6{,}398{,}744 & 56{,}241{,}592 \\
\bottomrule
\end{tabularx}}

\field{Size check:}{raw, $336{,}776\times19=6{,}398{,}744\ge3{,}000{,}000$, so the requirement is
already met. After encoding, $336{,}776\times167=56{,}241{,}592$. The 167 attributes are 28
numeric columns (12 from the flights file, with clock times in minutes; 8 joined weather values;
a \code{gusty} flag; and 7 derived columns: \code{gain}, \code{cancelled}, \code{wet},
\code{incomplete}, \code{weather\_missing}, \code{weekend}, \code{late15}), 4 sine/cosine columns
for departure clock time and day of year, and 135 one-hot columns: \code{dest} (105),
\code{carrier} (16), \code{weekday} (7), \code{season} (4) and \code{origin} (3).
\code{analysis.py} builds this matrix and checks its shape.}

\minihead{Key features}
{\footnotesize
\begin{tabularx}{\linewidth}{@{}L{2.2cm} L{1.5cm} Y@{}}
\toprule
\textbf{Feature} & \textbf{Type} & \textbf{Real-world meaning} \\
\midrule
\code{dep\_delay} & continuous & minutes the plane left the gate after schedule (negative = early) \\
\code{arr\_delay} & continuous & minutes late at the destination gate \\
\code{gain} & continuous & derived: \code{dep\_delay} $-$ \code{arr\_delay}, minutes made up en route \\
\code{carrier} & nominal & operating airline, 16 levels \\
\code{cancelled} & binary & derived: the flight never departed \\
\code{precip} & continuous & precipitation (inches) at the origin in the scheduled hour \\
\code{wet} & binary & derived: \code{precip} $>0$ \\
\code{season} & nominal & from month: Dec--Feb winter, Mar--May spring, Jun--Aug summer, Sep--Nov autumn \\
\code{origin}, \code{dest} & nominal & 3 NYC airports; 105 destinations \\
\bottomrule
\end{tabularx}}

%-----------------------------------------------------------------------
\section{Data preparation}

{\footnotesize
\begin{tabularx}{\linewidth}{@{}L{1.9cm} L{2.3cm} Y@{}}
\toprule
\textbf{Issue found} & \textbf{What we did} & \textbf{Why, and not the alternative} \\
\midrule
8{,}255 missing \code{dep\_time}, \code{dep\_delay}
 & coded \code{cancelled} $=1$; kept the rows
 & a missing departure is a cancellation. Dropping the rows would delete Test 4's response. \\
1{,}175 departed, no \code{arr\_delay}
 & flagged \code{incomplete}; left out of arrival tests only
 & diverted or unrecorded arrivals. An imputed arrival delay would create a made-up gain in
   Test 1. \\
2{,}512 missing \code{tailnum}
 & left as is
 & all 2{,}512 are cancelled flights (checked), so no aircraft was assigned \\
clock times as integers, \code{2400}
 & converted to minutes after midnight, mod 1440
 & 12:59 to 13:00 is 1 minute, not 41 \\
wind speed 1{,}048 mph (1 row)
 & set to missing
 & physically impossible; a sensor error \\
\code{wind\_gust} 79.6\% missing
 & replaced by a 0/1 \code{gusty} flag
 & gusts are recorded only when there is one, so missing means none \\
165 weather hours absent
 & left join on (origin, UTC hour); 1{,}556 flights flagged \code{weather\_\allowbreak missing}
 & an inner join would drop these flights from every test; only Test 2 leaves them out. The UTC
   hour avoids the repeated 1~a.m.\ hour when clocks go back (3 such rows). \\
27{,}789 arrival delays above the IQR fence (60.5 min)
 & kept, not capped
 & long delays are real. Capping would pull every mean down; the tests below are chosen to cope
   with the skew. \\
duplicates
 & checked full rows, (date, carrier, flight, origin, scheduled time), and (origin, hour) in
   weather
 & none found \\
\code{dest} 105, \code{carrier} 16 levels
 & one-hot; clock time and day of year as sine/cosine
 & codes have no order; sine/cosine keep 23:00 next to 00:00 \\
\bottomrule
\end{tabularx}}

%-----------------------------------------------------------------------
\section{What changed after preparation}

{\footnotesize
\begin{tabularx}{\linewidth}{@{}L{2.9cm} L{1.5cm} Y@{}}
\toprule
\textbf{Measure} & \textbf{Before} & \textbf{After} \\
\midrule
Rows & 336{,}776 & 336{,}776 (none dropped) \\
Attributes & 19 & 167 \\
Unexplained missing \code{dep\_delay} & 8{,}255 & 0: all labelled \code{cancelled} (2.45\%) \\
Unexplained missing \code{arr\_delay} & 9{,}430 & 0: 8{,}255 cancelled + 1{,}175 \code{incomplete} \\
Flights with weather & 0 (separate table) & 335{,}220; 1{,}556 flagged \\
Impossible wind values & 1 & 0 \\
Clock values \code{2400} & 179 & 0 \\
Duplicate rows & 0 & 0 \\
Arrival delays above IQR fence & 27{,}789 (8.5\%) & 27{,}789, kept on purpose \\
\bottomrule
\end{tabularx}}

\egblock{No row was lost. The 8{,}255 missing departures became a response variable
(\code{cancelled}). Arrival delay keeps its skewness of 3.7 and excess kurtosis of 29 because the
long tail was kept, so each test below is checked with a rank-based or day-level version.}

%-----------------------------------------------------------------------
\section{Basic questions and problem statement}

\begin{figure}[!ht]
  \figfile{fig_arrdelay.pdf}{3.6cm}
  \caption{57.7\% of completed flights arrive early (median $-5$ min), but the long right tail
  (skewness 3.7) pulls the mean up to $+6.9$.}
\end{figure}

\begin{figure}[!ht]
  \figfile{fig_gain.pdf}{3.6cm}
  \caption{Flights that leave on time make up about 6 min en route; flights that leave 30 min to
  2 h late make up less, not more.}
\end{figure}

\begin{figure}[!ht]
  \figfile{fig_hour_wet.pdf}{3.6cm}
  \caption{Delay builds through the day, and hours with precipitation sit above dry hours at every
  departure time.}
\end{figure}

\begin{figure}[!ht]
  \figfile{fig_cancel_daily.pdf}{3.6cm}
  \caption{The median day cancels 0.73\% of flights; on 8 and 9 February (a blizzard) more than
  half were cancelled.}
\end{figure}

\minihead{Basic questions}
\begin{enumerate}[label=\arabic*.]
  \item Which airline has the highest average arrival delay? Frontier (F9, 21.9 min), but it has
  only 681 completed flights; among the 11 large carriers, AirTran (FL, 20.1 min).
  \item Which month or season has the most cancellations? February (1{,}261 flights, 5.05\%);
  winter by season (3.50\%).
  \item Are departure and arrival delays related? Yes: $r=0.91$ over all completed flights and
  0.94 over flights that left 15 or more minutes late.
\end{enumerate}

\minihead{Our problem statements}

\noindent\fbox{\begin{minipage}{\dimexpr\linewidth-2\fboxsep-2\fboxrule\relax}
\textbf{P1 (primary).} Among flights that leave at least 15 minutes late, is the average time
made up in the air more than 5 minutes?

\smallskip
\textbf{P2.} Is a flight scheduled in an hour with precipitation at its origin delayed more than
15 minutes longer, on average, than one scheduled in a dry hour?

\smallskip
\textbf{P3.} Do carriers differ in arrival delay, and by 10 minutes or more for any pair?

\smallskip
\textbf{P4.} Does the cancellation rate depend on season, and does that survive removing the worst
storm days?

\smallskip
\textbf{P5.} Do weekend days have a different mean arrival delay from weekdays?
\end{minipage}}

\field{How they differ from the basic questions:}{the basic questions summarise one column each.
P1, P2 and P4 need columns that no raw row has (\code{gain}, \code{wet} from the hourly weather
join, \code{cancelled}), and each statement fixes a threshold for a difference that matters.}

%-----------------------------------------------------------------------
\section{Statistical inferences}

\field{Rows tested:}{Test 1: the 72{,}420 completed flights with \code{dep\_delay} $\ge15$.
Test 2: the 326{,}993 departed flights with weather (21{,}086 wet, 305{,}907 dry). Test 3: the
325{,}041 completed flights of the 11 carriers with at least 1{,}000 of them (AS, F9, YV, HA and
OO, 2{,}305 flights, are left out). Test 4: all 336{,}776 scheduled flights, summarised as 365
daily rates. Test 5: the 327{,}346 completed flights, summarised as 365 daily means.}

\field{How the hypotheses are set:}{P1 and P2 ask whether an effect is larger than a threshold,
so the threshold is the boundary of $H_0$. P3 to P5 ask whether groups differ, so $H_0$ is no
difference and the threshold is used to judge the confidence interval. Flights on the same day
share weather and congestion, so they are not independent: the main tests use standard errors
clustered by day (a t distribution with one less degree of freedom than the number of days) or
take the day as the unit.}

%------------------------------------------------------------ Test 1
\subsection*{Test 1 \textbullet\ In-flight delay recovery (P1): do late flights make up more than
5 minutes in the air?}

\tpart{a}{Hypotheses, written before running the test}
{\small
\begin{tabularx}{\linewidth}{@{}K{3.4cm} Y@{}}
\toprule
$H_0$ & $\mu_g\le5$ \\
$H_1$ & $\mu_g>5$ \\
Significance level & $\alpha=0.05$ \\
Side & one-sided (right) \\
Test & one-sample t-test on $g$ (paired design), standard errors clustered by day \\
Smallest difference that matters & 5 minutes \\
\bottomrule
\end{tabularx}}

\egblock{$g=\code{dep\_delay}-\code{arr\_delay}$ for one flight; $\mu_g$ is its mean over
flights that left at least 15 minutes late. 5 minutes is a third of the FAA's 15-minute line for a
delayed flight; a smaller gain rarely changes whether a passenger arrives late.}

\tpart{b}{Other tests we considered}
\egblock{Each flight gives two measurements of the same trip, so the data are paired and the test
works on the differences $g$. With $n=72{,}420$ the mean of $g$ is close to normal despite its
skewness ($-1.3$). Late flights on one day share the same congestion, so the standard error is
clustered by day (365 days).}
{\footnotesize
\begin{tabularx}{\linewidth}{@{}L{2.4cm} Y@{}}
\toprule
\textbf{Also considered} & \textbf{Why we did not use it} \\
\midrule
Paired t-test, flights independent & ignores the day clustering; its standard error is 7 times
smaller. Result: $t=-15.18$ (72{,}419), CI 3.63 to 3.95. \\
Wilcoxon signed-rank on $g-5$ & tests the pseudo-median, which differs from the mean when $g$ is
skewed. Result: $z=8.08$, $p<0.001$ in favour of $>5$. \\
Sign test against 5 & tests the median and ignores sizes: 38{,}125 flights gain more than 5 min
and 32{,}510 less, $p<0.001$. Median gain 7 min (day-bootstrap CI 6 to 7). \\
Welch two-sample t & treats departure and arrival delays as unrelated samples; against 0 it gives
$t=11.5$ instead of 47.5 for the paired test \\
z-test & needs a known $\sigma$ \\
\bottomrule
\end{tabularx}}

\tpart{c}{Result}
{\small
\begin{tabularx}{\linewidth}{@{}K{3.4cm} Y@{}}
\toprule
Estimate & $+3.79$ min (left 64.6 min late, arrived 60.9 min late) \\
95\% confidence interval & 2.71 to 4.87 min (day-clustered) \\
Test statistic (d.f.) & $t=-2.21$ (364) \\
p-value & 0.99 (one-sided) \\
Decision & fail to reject $H_0$ \\
\bottomrule
\end{tabularx}}

\egblock{Late flights make up about 4 minutes on average, 6\% of their 64.6-minute delay. Against
zero the gain is clear ($t=6.92$), but the whole interval lies below 5 minutes, and the reverse
test ($H_1{:}\ \mu_g<5$) gives $p=0.014$. Cohen's $d_z=0.18$.}

\egblock{The median tells a different story: the typical late flight gains 7 minutes, and the
Wilcoxon and sign tests put it above 5. The mean is lower because a minority of flights lose a
lot of time in the air. So the answer to P1, which is about the average, is no.}

\egblock{The gain is also not a catch-up. Flights that left less than 15 minutes late gain more
(6.2 min), the gain does not grow with departure delay (Spearman $\rho=-0.02$), and late flights
fly their route about as fast as the rest (air time 0.85 min above the route median, against
1.17 for other flights). It looks like schedule padding that every flight gets (Figure 2).}

\egblock{With one mean per day instead (365 days, each weighted equally) the gain is 5.86 min
(CI 4.99 to 6.72; $t=1.94$, $p=0.027$ against 5). This is a different average: on the quarter of
days with the fewest late flights they gain 10.1 min, on the quarter with the most they gain
nothing ($-0.1$ min; Spearman $\rho=-0.46$ between the two). Weighting days equally gives the
quiet days more say. On a bad day, late flights do not make up time.}

%------------------------------------------------------------ Test 2
\subsection*{Test 2 \textbullet\ Weather impact on departures (P2): does precipitation add more
than 15 minutes of delay?}

\tpart{a}{Hypotheses, written before running the test}
{\small
\begin{tabularx}{\linewidth}{@{}K{3.4cm} Y@{}}
\toprule
$H_0$ & $\mu_{\text{wet}}-\mu_{\text{dry}}\le15$ \\
$H_1$ & $\mu_{\text{wet}}-\mu_{\text{dry}}>15$ \\
Significance level & $\alpha=0.05$ \\
Side & one-sided (right) \\
Test & difference in means (Welch form), standard errors clustered by day \\
Smallest difference that matters & 15 minutes \\
\bottomrule
\end{tabularx}}

\egblock{$\mu_{\text{wet}}$, $\mu_{\text{dry}}$: mean departure delay of flights scheduled in an
hour with and without precipitation at their origin. 15 minutes is the FAA's line for a delayed
flight.}

\tpart{b}{Other tests we considered}
\egblock{Two independent groups, a continuous response and unknown $\sigma$ call for a two-sample
t-test; the spreads differ (SD 59.2 vs 38.3 min; Levene $W=3{,}916$, $p<0.001$), so the Welch
form. Flights in one hour share its weather, and the hourly mean delay at an airport is strongly
autocorrelated (lag-1 $r=0.81$), so standard errors are clustered by day.}
{\footnotesize
\begin{tabularx}{\linewidth}{@{}L{2.4cm} Y@{}}
\toprule
\textbf{Also considered} & \textbf{Why we did not use it} \\
\midrule
Welch t-test, flights independent & ignores that all flights in a wet hour share the same weather.
Result: $t=11.0$ (22{,}314), CI 18.7 to 20.4. \\
Pooled t-test & equal variances are rejected; against 0 it gives $t=68.8$ \\
Mann--Whitney U & compares ranks, not minutes. $P(\text{wet}>\text{dry})=0.64$, $p<0.001$;
medians 5 vs $-2$ min. \\
Permutation test & shuffling single flights breaks the clustering; gives only a p-value \\
z-test & needs a known $\sigma$ \\
\bottomrule
\end{tabularx}}

\tpart{c}{Result}
{\small
\begin{tabularx}{\linewidth}{@{}K{3.4cm} Y@{}}
\toprule
Estimate & $+19.6$ min (30.9 wet, 11.4 dry) \\
95\% confidence interval & 14.8 to 24.3 min (day-clustered) \\
Test statistic (d.f.) & $t=1.88$ (363) \\
p-value & 0.030 (one-sided) \\
Decision & reject $H_0$ at 0.05 \\
\bottomrule
\end{tabularx}}

\egblock{Precipitation goes with about 20 more minutes of departure delay, and against zero this is
clear ($t=8.09$). That it is more than 15 minutes is borderline. The 95\% interval starts at 14.8;
the one-sided test at 0.05 matches the 90\% interval, 15.6 to 23.5, which just clears 15. 40\% of
wet-hour departures are 15 or more minutes late against 21\% of dry ones, and delay rises with
intensity: 27.3, 30.4 and 40.9 min for up to 0.02, 0.02 to 0.1 and over 0.1 inch.}

\egblock{Two checks. (i) One row per airport-hour (19{,}326 hours) gives $+20.5$ min (CI 18.3 to
22.6; Welch $t=5.04$, d.f.\ 1{,}328, $p<0.001$ against 15), but neighbouring hours are still
correlated. (ii) Rain may fall more often in months and hours that have more delay anyway. With
month, hour and airport fixed effects (errors clustered by day) the gap is $+18.5$ min (CI 14.1 to
22.9; $t=1.57$, $p=0.059$), and $H_0$ is no longer rejected. Precipitation clearly adds delay;
that it adds more than 15 minutes is not established.}

%------------------------------------------------------------ Test 3
\subsection*{Test 3 \textbullet\ Carrier performance (P3): do carriers differ in arrival delay?}

\tpart{a}{Hypotheses, written before running the test}
{\small
\begin{tabularx}{\linewidth}{@{}K{3.4cm} Y@{}}
\toprule
$H_0$ & the arrival-delay distributions of the 11 carriers are identical \\
$H_1$ & at least one carrier's delays tend to be larger or smaller \\
Significance level & $\alpha=0.05$ \\
Side & omnibus (upper tail of $H$) \\
Test & Kruskal--Wallis, then Dunn's test with Holm adjustment \\
Smallest difference that matters & 10 minutes between two carriers \\
\bottomrule
\end{tabularx}}

\egblock{10 minutes is two-thirds of the FAA line, a gap a frequent traveller would notice. A rank
test cannot put minutes in $H_0$, so the threshold is applied to simultaneous 95\% intervals for
the pairwise mean differences (Games--Howell).}

\tpart{b}{Other tests we considered}
\egblock{Eleven independent groups, a heavily skewed response (per-carrier skewness 2.6 to 5.4)
and unequal spreads (Levene $W=207.5$, $p<0.001$). Kruskal--Wallis compares ranks, so a few
20-hour delays cannot dominate it, and it does not need normality. Dunn's test is the matching
rank-based follow-up.}
{\footnotesize
\begin{tabularx}{\linewidth}{@{}L{2.4cm} Y@{}}
\toprule
\textbf{Also considered} & \textbf{Why we did not use it} \\
\midrule
One-way ANOVA & assumes equal variances, which Levene rejects. Result: $F=477.7$ (10, 325{,}030),
same decision. \\
Welch ANOVA & handles unequal variances but compares means, which the tails drive. Result:
$F=445.5$ (10, 50{,}794), $\eta^2=0.014$. \\
55 pairwise t-tests & the chance of at least one false rejection would be $1-0.95^{55}=94\%$ \\
\bottomrule
\end{tabularx}}

\tpart{c}{Result}
{\small
\begin{tabularx}{\linewidth}{@{}K{3.4cm} Y@{}}
\toprule
Estimate & worst (AirTran, FL) minus best (American, AA): $+19.8$ min mean (20.1 vs 0.4);
medians 5 vs $-9$ \\
95\% confidence interval & 16.6 to 22.9 min (Games--Howell, simultaneous); median gap 13 to 16
(bootstrap) \\
Test statistic (d.f.) & $H=8{,}170.2$ (10) \\
p-value & $<0.001$ \\
Decision & reject $H_0$ \\
\bottomrule
\end{tabularx}}

\begin{figure}[!ht]
  \figfile{fig_carriers.pdf}{3.6cm}
  \caption{Mean arrival delay by carrier with day-clustered 95\% CIs. The grey band spans 10 min
  above the best carrier; only EV and FL sit clearly outside it.}
\end{figure}

\egblock{Carriers differ, but carrier explains little: $\varepsilon^2=0.025$, about 2.5\% of the
variation in ranks. Dunn--Holm finds 53 of 55 pairs different (not AA--VX or B6--WN). Only 14
pairs differ by 10 minutes or more in mean, and 13 of them involve ExpressJet (EV, 15.8 min) or
AirTran (FL, 20.1); the other is MQ against AA. Using the simultaneous intervals, 10 pairs are
more than 10 minutes apart for certain, 34 are less than 10 apart for certain, and 11 are
undecided. By medians only 6 pairs reach 10 minutes, all involving FL.}

\egblock{Because flights on one day share conditions, we checked the result with the day built
in. A Friedman test with days as blocks (each carrier's median delay per day, 365 days) gives
$\chi^2=1{,}144.6$ (10), $p<0.001$, Kendall's $W=0.31$, so the order holds within days. For
FL $-$ AA, day-clustered errors give CI 17.1 to 22.4 and day fixed effects give $+20.1$ (CI 17.4
to 22.9). Kruskal--Wallis within each airport is significant at all three (EWR $H=3{,}306$, 9
d.f.), so the gap is not about which airport a carrier uses.}

%------------------------------------------------------------ Test 4
\subsection*{Test 4 \textbullet\ Seasonal cancellation risk (P4): do cancellations depend on
season?}

\tpart{a}{Hypotheses, written before running the test}
{\small
\begin{tabularx}{\linewidth}{@{}K{3.4cm} Y@{}}
\toprule
$H_0$ & the daily cancellation rate has the same distribution in all four seasons \\
$H_1$ & it differs in at least one season \\
Significance level & $\alpha=0.05$ \\
Side & omnibus (upper tail of $H$) \\
Test & Kruskal--Wallis on the 365 daily cancellation rates \\
Smallest difference that matters & 1 percentage point between two seasons \\
\bottomrule
\end{tabularx}}

\egblock{Against a yearly rate of 2.45\%, a 1-point gap is a 40\% change in a traveller's risk,
about 840 extra cancelled flights per season out of New York.}

\tpart{b}{Other tests we considered}
\egblock{Cancellations come in bursts: the 10 worst days hold 27.2\% of the year's cancellations
(Figure 4). Flights on one day are cancelled together, so the day is the unit. Daily rates are
very skewed (median 0.73\%, maximum 57\%) and their spread differs by season (Levene $W=3.1$,
$p=0.026$), so a rank test.}
{\footnotesize
\begin{tabularx}{\linewidth}{@{}L{2.4cm} Y@{}}
\toprule
\textbf{Also considered} & \textbf{Why we did not use it} \\
\midrule
Chi-square test of independence ($4\times2$) & counts are large enough (smallest expected
1{,}963), but it treats 336{,}776 flights as independent. Result: $\chi^2=1{,}055.3$ (3); its CI
for winter $-$ autumn (Newcombe, 2.26 to 2.55 pp) is 12 times narrower than the day-clustered
one. \\
G-test, Fisher, logistic regression & same independence assumption as chi-square ($G=1{,}173$) \\
ANOVA on daily rates & daily means are driven by two storm days. Welch form: $F=7.26$ (3, 190),
$p<0.001$. \\
\bottomrule
\end{tabularx}}

\tpart{c}{Result}
{\small
\begin{tabularx}{\linewidth}{@{}K{3.4cm} Y@{}}
\toprule
Estimate & rates: winter 3.50\%, summer 2.80\%, spring 2.43\%, autumn 1.10\%; winter $-$ autumn
$=2.40$ pp \\
95\% confidence interval & 0.69 to 4.12 pp (day-clustered) \\
Test statistic (d.f.) & $H=39.9$ (3) \\
p-value & $<0.001$ \\
Decision & reject $H_0$ \\
\bottomrule
\end{tabularx}}

\egblock{Seasons differ (a day-clustered Wald test on the flight-level rates agrees: $F=6.99$ (3,
364), $p<0.001$). But the interval for the largest gap reaches down to 0.69, below the 1-point
threshold. Pairwise day-clustered tests (Holm) separate only autumn from each of the other three;
winter, spring and summer cannot be told apart.}

\egblock{Winter's lead comes from one storm: 472 flights were cancelled on 8 February and 393 on 9
February. Dropping those two days puts winter (2.47\%) below summer (2.80\%). Dropping the 10
worst days, the test still rejects ($H=39.7$, $p<0.001$), with summer 2.69\%, winter 1.92\%,
spring 1.82\% and autumn 0.88\%; winter $-$ autumn is then 1.04 pp (CI 0.43 to 1.64). On a typical
day (median daily rate) the gap is 0.67 pp (winter 1.08\%, autumn 0.41\%). Autumn is the lowest
season in every version; that any season differs by more than 1 point is not shown.}

%------------------------------------------------------------ Test 5
\subsection*{Test 5 \textbullet\ Weekly delay patterns (P5): weekend versus weekday arrival
delay}

\tpart{a}{Hypotheses, written before running the test}
{\small
\begin{tabularx}{\linewidth}{@{}K{3.4cm} Y@{}}
\toprule
$H_0$ & $\mu_{\text{we}}-\mu_{\text{wd}}=0$ \\
$H_1$ & $\mu_{\text{we}}-\mu_{\text{wd}}\ne0$ \\
Significance level & $\alpha=0.05$ \\
Side & two-sided \\
Test & Welch t-test on daily mean arrival delay (104 weekend days, 261 weekdays) \\
Smallest difference that matters & 5 minutes \\
\bottomrule
\end{tabularx}}

\egblock{$\mu_{\text{we}}$, $\mu_{\text{wd}}$: mean daily arrival delay on Saturday--Sunday and
Monday--Friday. Two-sided, because fewer flights (less congestion) and more leisure travellers
(slower boarding) push in opposite directions. Weekend is a property of a day, so the day is the
unit.}

\tpart{b}{Other tests we considered}
\egblock{Two independent groups of days, unknown $\sigma$, and unequal spreads of the daily means
(SD 13.4 vs 18.2 min; Levene $W=6.0$, $p=0.015$), so the Welch form.}
{\footnotesize
\begin{tabularx}{\linewidth}{@{}L{2.4cm} Y@{}}
\toprule
\textbf{Also considered} & \textbf{Why we did not use it} \\
\midrule
Welch t-test on flights & treats flights on one day as independent: $t=-39.1$ (160{,}784), CI
$-6.9$ to $-6.3$. Clustered by day: $-6.6$ min (CI $-9.9$ to $-3.3$). \\
Pooled t-test & equal variances are rejected; $t=-3.66$ on days \\
Mann--Whitney U & compares ranks, not minutes. On daily means $p<0.001$; on flights
$P(\text{we}>\text{wd})=0.45$, medians $-8$ vs $-4$ min. \\
Paired t-test & a day belongs to one group only \\
\bottomrule
\end{tabularx}}

\tpart{c}{Result}
{\small
\begin{tabularx}{\linewidth}{@{}K{3.4cm} Y@{}}
\toprule
Estimate & $-7.2$ min (1.85 weekend, 9.06 weekday) \\
95\% confidence interval & $-10.6$ to $-3.8$ min \\
Test statistic (d.f.) & $t=-4.15$ (254.8) \\
p-value & $<0.001$ \\
Decision & reject $H_0$ \\
\bottomrule
\end{tabularx}}

\begin{figure}[!ht]
  \figfile{fig_weekday.pdf}{3.6cm}
  \caption{Daily mean arrival delay by day of the week. Saturday is lowest; Sunday looks like a
  weekday.}
\end{figure}

\egblock{Weekend days have lower arrival delay. Whether the gap is more than 5 minutes is
uncertain, because the interval reaches $-3.8$. Most of it is Saturday, the day with the fewest
flights (38{,}720 scheduled, against about 50{,}000 on each weekday): Saturday flights average
$-1.45$ min, 10.0 min less than weekday flights (CI $-13.6$ to $-6.4$, day-clustered). Sunday
averages 4.82, 3.75 less (CI $-8.1$ to $+0.6$), which cannot be told apart from a weekday.
Thursday is the worst day at 11.74 min.}

\minihead{Summary of the five tests}
{\footnotesize
\begin{tabularx}{\linewidth}{@{}L{0.3cm} L{2.4cm} L{1.4cm} Y@{}}
\toprule
\textbf{\#} & \textbf{Estimate (95\% CI)} & \textbf{Decision} & \textbf{Matters?} \\
\midrule
1 & gain $+3.8$ min (2.7, 4.9) & fail to reject & no: mean below 5 min (median 7); padding, not
catch-up \\
2 & wet $+19.6$ min (14.8, 24.3) & reject ($p=0.030$) & borderline: $p=0.059$ after calendar
adjustment \\
3 & FL $-$ AA $+19.8$ min (16.6, 22.9) & reject & only pairs with EV or FL; $\varepsilon^2=0.025$ \\
4 & winter $-$ autumn $+2.4$ pp (0.7, 4.1) & reject & not shown above 1 pp; only autumn stands
apart \\
5 & weekend $-7.2$ min ($-10.6$, $-3.8$) & reject & real; above 5 min uncertain; a Saturday
effect \\
\bottomrule
\end{tabularx}}

\egblock{With five tests, a Bonferroni level of $0.05/5=0.01$ changes one decision: Test 2 would
not reject. With 300{,}000 rows almost every flight-level test rejects, so the useful columns are
whether the interval clears the threshold and whether the result holds with the day as the unit.}

%-----------------------------------------------------------------------
\section{Limitations}

\begin{itemize}
  \item \textbf{Association, not cause.} Precipitation comes with low visibility, wind and
  congestion; Test 2 cannot separate them, and its fixed-effects check removes only calendar and
  airport effects.
  \item \textbf{Dependence.} Clustering by day handles flights that share a day, but consecutive
  days are also correlated (lag-1 $r=0.37$ for daily mean delay), so the day-level p-values are
  still somewhat optimistic. The Kruskal--Wallis p-value in Test 3 is flight-level; the Friedman
  and clustered checks cover it.
  \item \textbf{All flights, not a sample.} The data are every 2013 NYC departure, so p-values
  measure how unusual a pattern is against chance, not how well it carries to other years.
  \item \textbf{Survivor bias.} Delay tests use flights that flew. The worst hours cancel flights
  instead of delaying them, so Test 2 likely understates the effect of bad weather.
  \item \textbf{Weather at the origin, at the scheduled hour only.} Storms at the destination or a
  few hours earlier are not measured.
  \item \textbf{Choices we made.} The 5, 10 and 15-minute and 1-point thresholds follow from the
  FAA's 15-minute rule and the base rate; others would change the Matters column. Winter is
  January, February and December of 2013, so it joins two different winters.
  \item \textbf{Scope.} One year, one city, domestic flights. AirTran has since merged into
  Southwest and US Airways into American, so Test 3 is a 2013 snapshot.
\end{itemize}

\end{document}
